\section{近期研究（2026 年 8--11 月）}
\paragraph{3DGSI-Assessor：区分几何和颜色的压缩质量评测}
2026.08~\cite{2608_03279}。常用的 PSNR 和 SSIM 将视图的差异压成一个分数，难以判断高斯压缩究竟损伤了物体边界还是颜色。3DGS-IEval-15K+ 为十个真实场景的 15200 张图像分别收集总体、几何和颜色评分，共形成 45600 个平均主观分。3DGSI-Assessor 用视觉语言对齐和 LoRA 微调，使同一个模型给出三个维度的语言等级和数值分数，便于区分不同失真。

论文在该库上按每个维度对照主观分，比较 SRCC、PLCC 和 KRCC。作者另外将各学习模型在该库训练后直接迁移到 ENeRF-QA、NeRF-VSQA 和 GSC-QA，通过全数据集相关性检验跨域迁移。总体相关性与按场景作相对排序的协议不同。数据仍限于这十个真实场景及所覆盖的压缩失真，评测设计本身也未验证将分数直接作为率失真训练目标的收益。

\paragraph{用多视图预测置信度约束高斯表面}
2026.08~\cite{2608_06117}。逐视图预测的深度和法向可以约束 3DGS，但不同相机下的预测在反射表面尤其容易不一致。该研究比较单目几何先验与 VGGT 的多视图深度和法向先验，再用 VGGT 预测的置信图对几何正则加权。在可靠区域，先验帮助光度拟合选定表面；在低置信区域，误预测对表面拟合的影响被减轻。

论文在 DTU、Shiny Blender 和 Tanks and Temples 对照几何质量，主要增益出现在 Shiny Blender 的反射物体；DTU 和 Tanks and Temples 的整体改进较小。Shiny Blender 官方测试视角的渲染 PSNR 与 PGSR 相近。这组结果支持把置信度用作几何先验的调节量，同时表明几何误差下降不必然体现在图像指标上。

\paragraph{TRACE-GS：在稀疏去噪轨迹上蒸馏稠密几何}
2026.08~\cite{2608_10286}。少视角训练出的高斯在远离输入相机的方向容易产生缺损。TRACE-GS 把视图恢复模型训练成稀疏输入的修复器，训练时为教师提供由额外视图建立的稠密几何。学生沿自己的扩散去噪轨迹访问潜变量状态，教师在同一状态提供去噪方向和跨视图查询监督，不需要单独运行教师轨迹。部署仅使用稀疏条件学生，并将修复后的视图用于迭代更新原始高斯。

在与训练场景不重叠的 28 个 DL3DV 基准场景上，作者按 9、6、3 幅输入评估，相对原始 3DGS 报告的 PSNR 增益分别为 2.02、2.35 和 3.20 dB。论文还在 Mip-NeRF 360 和 NeRFBusters 验证域外表现。基准比较应注明相同的目标视图和重建协议。方法针对静态场景，当前部署仍包含视频扩散恢复与交替高斯优化。

\paragraph{GS-CPE：从场景检索到可见性约束的位姿精修}
2026.08~\cite{2608_10938}。GS-CPE 接收查询图像和已训练的三维高斯场景。粗定位先检索与查询相似的数据库视图，再用高斯场景提供的几何线索迭代建立二维至三维对应。精定位用可见性掩码排除遮挡区域，对查询图像与当前姿态生成的视图作多尺度变形对齐，并在优化过程中自适应重渲染。高斯场景在这里充当定位地图与可渲染参考。

7Scenes 上按场景汇总的平移和旋转中位误差由粗阶段的 1.08 厘米、0.31 度降低到 0.91 厘米、0.29 度；Cambridge Landmarks 上由 23 厘米、0.44 度降到 21 厘米、0.39 度。FAST-LIVO2 的 CBD2 序列仍有基线更准，论文将差距与初始化误差导致的场景渲染质量联系起来。定位结果依赖预建场景，该原文用基准参考位姿或由重建得到的位姿评价，并未在所有数据集上使用独立测量真值。

\paragraph{M³ISR：把角度变化、码率与 streaming 成本纳入同一基准}
2026.08~\cite{2608_22465}。M³ISR 用受控合成采集比较静态和动态高斯方法，提供 25 个室内外场景。六路同步相机覆盖 120 度水平扇区，共用观测中心以压低平移视差，使模型的角度外推更容易单独评价。数据还提供校准参数和稠密标注，挑战涵盖静态与动态的新视图合成、流式表示及压缩。

论文给出部分赛道的代表性基线，并要求压缩对照把解码所需的参数、潜变量、先验和元数据一并计入码率。现有压缩参考曲线的视图协议等细节并不完整，因此不宜据此得出某种前馈编码器已占优的结论。两秒的合成序列和有限相机角度也不能用于判断真实环境下的长期漂移、传感器噪声影响或完整 360 度覆盖。

\paragraph{PIVOT：区分已见轨迹插值和整条新轨迹外推}
2026.08~\cite{2608_25401}。常见的新视图协议在训练轨迹内留出若干帧，容易把同一采集分布内的插值和新轨迹上的外推混为一谈。PIVOT 在五个真实场景保存轨迹身份、测量位姿与优化位姿、标定内参与优化内参，使这些因素可以分别改变。作者以从评估轨迹指向训练位姿集合的有向位姿 Chamfer 距离描述覆盖差异。

在 Nerfacto 与 Splatfacto 两个实现的初步评估中，五个场景的整条未见轨迹在 PSNR、SSIM 和 LPIPS 上都比已见轨迹内留出帧更难。位姿距离相近的轨迹也可能有不同渲染误差，因此距离适合用来描述协议，不能替代对遮挡和场景内容的分析。这一结果的范围是所测五个场景与两种实现。

\paragraph{把水下姿态可解性与渲染质量分开测量}
2026.08~\cite{2608_25483}。水下高斯方法可以在图像端预处理介质效应、单独学习介质场，或将介质效应附着在高斯原语上。该研究将这些设计与原始 3DGS 放在同一协议中比较。实验分别固定轨迹、位姿和初始化后改变浑浊度，以及在清水条件下将编码器位姿替换为 COLMAP 位姿，从而把渲染器退化与上游注册失败分开。

论文报告实地测量场景中，固定图像恢复预处理加原始 3DGS 得到的表面优于所测介质模型。浑浊度升高时，位姿注册可能先于渲染器失效；相机随动灯光也削弱了所测介质模型的优势。这些结论限定在所测四组水域条件和比较对象，比较时不能只看渲染外观。

\paragraph{View-CP：将高斯置信区间校准到整幅视图}
2026.09~\cite{2609_10307}。把不同视图的像素混在一起校准，即使总体平均覆盖率正确，也可能让某些视图的大量像素同时落在区间之外。View-CP 目标是让新视图达到给定比例的像素颜色覆盖。它将渲染器提供的像素空间误差形状与一个视图难度因子分开，再用留出视图上的比值分位数校准 RGB 预测盒。

13 个真实场景的对照中，像素混合基线的像素覆盖率为 89.9\%，但达到覆盖要求的视图仅占 61.4\%；目标概率为 90\% 时，View-CP 按不同尺度图得到 91.7\% 至 92.0\% 的视图事件覆盖率。保证针对可交换视图的边际概率，不能解释为每个指定相机位姿的条件保证。沿时间排序的轨迹可能破坏该假设，方法还需要有真值的目标场景留出视图。

\paragraph{从高斯原语失真出发设计评测协议}
2026.09~\cite{2609_23027}。相同的图像伪影可能来自高斯几何、属性量化、深度排序、滤波或轨迹采样。该综述将失真机制与可见症状对应起来，并区分需要参考视图的指标、直接用渲染视图判断的指标、继承自点云或网格的方法，以及直接读取高斯属性的方法。

作者引用受控评测说明，传统视图指标在某些高斯特有失真下与主观评分相关性不足。这一比较来自所综述的文献。对于新协议，文章建议同时保留规范视图、留出视图、针对特定失效的压力视图和短轨迹，并明确记录渲染器与压缩配置。该建议旨在区分质量损失的来源，文中没有提出并验证新的统一指标。

\paragraph{3DGS-SC：把屏幕文字可读性列为单独评测目标}
2026.09~\cite{2609_31756}。屏幕内容具有密集文字和规则边缘，PSNR 接近的渲染也可能相差很大的可读性。3DGS-SC 在程序化静态场景中保存屏幕与文字的精确标注。除了整图和屏幕区域的失真，协议单独评价 OCR 和带容差的边缘 F1，并在同一配置下比较五个代表性高斯方法。

LightGaussian 的屏幕 PSNR 为 20.71 dB、OCR 准确率为 73.5\%，Mip-Splatting 则分别为 20.63 dB 和 86.2\%。两者仅差 0.08 dB，文字识别却反向相差 12.7 个百分点。此例说明屏幕指标需要按任务选用。程序化基准还需要真实显示器采集和 OCR 参考校准来检验适用范围。

\paragraph{用立体主观实验揭示单目指标忽略的高斯伪影}
2026.09~\cite{2609_38525}。研究在减少采集视图后，用两种初始化训练其余设置相同的高斯场，再比较单目或立体显示的偏好。设计使评价更集中在初始化导致的结构差异，而不把它与渲染器或训练配置差异混在一起。

20 位参与者在立体条件共作出 560 次选择，其中 438 次偏好联合初始化，为 78.2\%；单目条件的对应汇总为 58.4\%。自动指标的差异较小且不一致。HMD 中呈现的是静态立体图像对，结果不覆盖自由六自由度探索、运动中的时间伪影或长时舒适度。

\paragraph{SPLATIFY：用可训练的高斯组件约束论文复现}
2026.10~\cite{2610_09116}。SPLATIFY 将论文中的改动限制在 gsplat 方法模板的扩展位点，例如损失、增密和优化模块。系统从引用中定位候选组件，按依赖关系组合和生成实现，再通过训练日志、渲染结果与原语诊断迭代。约束模板减少了生成代码破坏高斯训练流程的机会。

论文在 Splatify-Bench 对比专家和通用论文转代码系统，报告若干可训练实现的 PSNR 与专家实现相差不超过 0.5 dB，并在组合改进中报告最高 2.4 dB 的增益。这些是特定基准与比较流程的结果。渲染分数接近不能证明函数语义、训练轨迹和作者实现相同，审计仍须保留关键组件的出处与实测结果。

\paragraph{从端到端计时边界看实时高斯重建}
2026.10~\cite{2610_17513}。标称实时的高斯渲染器可能只说明已有场景的显示速度。该综述按重建管线组织先前研究，将前馈生成、动态表示、参数压缩、流式更新与 SLAM 放到各自的时延边界下讨论。例如，动态参数化涉及时间稳定性，剪枝、量化与细节层级涉及体积和画质，光栅化选择又改变伪影与合成成本。

文中对已有实验要求同时记录输入设置、分辨率、硬件和计时边界。前馈重建的冷启动、每帧地图更新和已有模型的渲染帧率对应不同任务；没有统一这些条件的结果，只能用来理解各个方法的表现。这些内容是综述提出的分析协议。

\paragraph{SplatXR：将高斯渲染、几何和交互集成到头显工作流}
2026.10~\cite{2610_26338}。SplatXR 在 Unity 里把高斯画面与真实场景和可交互数字物体组织成多个层。工具链含头显端实时渲染、高斯转网格、物理交互和层间混合，也支持把头显拍摄的物体交给外部设备重建。设备端显示性能和外部重建成本需要分开统计。

作者在 Quest 3 上测试 9 个场景，每个取 5 万、10 万和 20 万三档高斯数，共 27 种渲染配置。网格生成与导入另在笔记本的 Unity 编辑器中测时。作者讨论的细节层级支持更高的静态高斯数，但静态模型序列对连续动态场景的表达和效率有限。

\paragraph{SplatLab：把高斯研究候选与实验记录对应起来}
2026.10~\cite{2610_30468}。SplatLab 将自动研究定位在已复现的高斯基线上。系统先定位论文、代码和数据，再提出修改候选。在多个相对独立的搜索群组中，候选根据实验结果排序与迭代，最后由审计阶段检查文稿陈述是否能对应到已记录的运行。这使结果可以追溯到某一组具体配置，而不是只有生成出的实验描述。

论文区分了以原始 3DGS 为对照的星云、气泡和钻石场景，以及在 Mip-NeRF 360 等已有基准上修改已发表方法的场景。作者报告的 0.18 至 3.78 dB 改进相对的是各自复现目标，不能看成同一基准上的统一增益。该过程还承担复现基础设施与搜索计算成本，PSNR 导向的结果需要结合其他质量指标和独立复现判断。
